\section{费米液体的输运与集体模}

Gross 第 32 章：由 Landau 动力学方程导出零声，并讨论碰撞导致的输运系数。

\subsection{无碰撞动力学}

准粒子分布 $n_{\mathbf{p}}(\mathbf{r},t)$ 的守恒形式（无碰撞）为
\begin{equation}
  \partial_t\delta n_{\mathbf{p}}
  +\mathbf{v}_{\mathbf{p}}\cdot\nabla_{\mathbf{r}}
  \bigl(
  \delta n_{\mathbf{p}}
  -\frac{\partial n^0}{\partial\varepsilon_{\mathbf{p}}}
  \delta\varepsilon_{\mathbf{p}}
  \bigr)
  =0,
\label{eq:landau_kinetic}
\end{equation}
其中 $\delta\varepsilon_{\mathbf{p}}=\sum_{\mathbf{p}'}f_{\mathbf{p}\mathbf{p}'}\delta n_{\mathbf{p}'}$，$\mathbf{v}_{\mathbf{p}}=\nabla_{\mathbf{p}}\varepsilon_{\mathbf{p}}$。第二项中减去 $(\partial n^0/\partial\varepsilon)\delta\varepsilon$ 是因为局域平衡已按重整化能量重新填充。

\subsection{零声色散}

设 $\delta n\propto e^{i(\mathbf{q}\cdot\mathbf{r}-\omega t)}$，并只保留 $F_0^s$ 各向同性形变。令
\begin{equation}
  \nu(\hat{\mathbf{p}})
  =\delta n_{\mathbf{p}}\Big/\Bigl(-\frac{\partial n^0}{\partial\varepsilon}\Bigr),
  \qquad
  s=\frac{\omega}{q v_F}.
\end{equation}
对费米面角积分后得
\begin{equation}
  \frac{s}{2}\ln\frac{s+1}{s-1}-1
  =\frac{1}{F_0^s}
  \qquad(F_0^s>0).
\label{eq:zerosound}
\end{equation}
$s>1$ 的根位于粒子--空穴连续谱 $\omega<q v_F$ 之上，无朗道阻尼，称为\textbf{零声}（碰撞可忽略的高频/短波极限，$\omega\tau\gg1$）。$F_0^s\to0$ 时 $s\to1^+$，模式并入连续谱边缘。

\subsection{第一声与碰撞}

当 $\omega\tau\ll1$，碰撞积分强迫局部热平衡，恢复流体力学：第一声速由压缩率决定，
\begin{equation}
  c_1^2
  =\frac{v_F^2}{3}
  \frac{1+F_0^s}{1+F_1^s/3}.
\end{equation}
黏滞 $\eta\sim p_F^2\tau_\eta/m^*$、热导 $\kappa_T\sim C_V v_F^2\tau_T$。费米液体相空间给出 $\tau\sim T^{-2}$，故纯净极限 $\rho\sim T^2$。

\subsection{与微观响应}

$F_\ell$ 由微观 $\Gamma^\omega$ 投影得到；零声极点应与密度响应 $\chi(q,\omega)$ 在相应通道的集体极点一致。这把第 5 章 RPA/局域场与第 7 章 Landau 参数连成一体。
